\documentclass[10pt,a4paper]{amsart}
\usepackage{amsmath,amssymb,amsthm,graphicx}
\usepackage{verbatim}
\newenvironment{reference}{\comment}{\endcomment}
%\newenvironment{reference}{}{}
\newenvironment{slogan}{\comment}{\endcomment}
\newenvironment{history}{\comment}{\endcomment}
% For commutative diagrams we use Xy-pic
\usepackage[all]{xy}

% We use 2cell for 2-commutative diagrams.
\xyoption{2cell}
% Theorem environments.
%
\theoremstyle{plain}
\newtheorem{theorem}[subsection]{Theorem}
\newtheorem{proposition}[subsection]{Proposition}
\newtheorem{lemma}[subsection]{Lemma}

\theoremstyle{definition}
\newtheorem{definition}[subsection]{Definition}
\newtheorem{example}[subsection]{Example}
\newtheorem{exercise}[subsection]{Exercise}
\newtheorem{situation}[subsection]{Situation}

\theoremstyle{remark}
\newtheorem{remark}[subsection]{Remark}
\newtheorem{remarks}[subsection]{Remarks}

\numberwithin{equation}{subsection}

% Macros
%
\def\lim{\mathop{\mathrm{lim}}\nolimits}
\def\colim{\mathop{\mathrm{colim}}\nolimits}
\def\Spec{\mathop{\mathrm{Spec}}}
\def\Hom{\mathop{\mathrm{Hom}}\nolimits}
\def\Ext{\mathop{\mathrm{Ext}}\nolimits}
\def\SheafHom{\mathop{\mathcal{H}\!\mathit{om}}\nolimits}
\def\SheafExt{\mathop{\mathcal{E}\!\mathit{xt}}\nolimits}
\def\Sch{\mathit{Sch}}
\def\Mor{\mathop{\mathrm{Mor}}\nolimits}
\def\Ob{\mathop{\mathrm{Ob}}\nolimits}
\def\Sh{\mathop{\mathit{Sh}}\nolimits}
\def\NL{\mathop{N\!L}\nolimits}
\def\CH{\mathop{\mathrm{CH}}\nolimits}
\def\proetale{{pro\text{-}\acute{e}tale}}
\def\etale{{\acute{e}tale}}
\def\QCoh{\mathit{QCoh}}
\def\Ker{\mathop{\mathrm{Ker}}}
\def\Im{\mathop{\mathrm{Im}}}
\def\Coker{\mathop{\mathrm{Coker}}}
\def\Coim{\mathop{\mathrm{Coim}}}

% Boxtimes
%
\DeclareMathSymbol{\boxtimes}{\mathbin}{AMSa}{"02}

%
% Macros for moduli stacks/spaces
%
\def\QCohstack{\mathcal{QC}\!\mathit{oh}}
\def\Cohstack{\mathcal{C}\!\mathit{oh}}
\def\Spacesstack{\mathcal{S}\!\mathit{paces}}
\def\Quotfunctor{\mathrm{Quot}}
\def\Hilbfunctor{\mathrm{Hilb}}
\def\Curvesstack{\mathcal{C}\!\mathit{urves}}
\def\Polarizedstack{\mathcal{P}\!\mathit{olarized}}
\def\Complexesstack{\mathcal{C}\!\mathit{omplexes}}
% \Pic is the operator that assigns to X its picard group, usage \Pic(X)
% \Picardstack_{X/B} denotes the Picard stack of X over B
% \Picardfunctor_{X/B} denotes the Picard functor of X over B
\def\Pic{\mathop{\mathrm{Pic}}\nolimits}
\def\Picardstack{\mathcal{P}\!\mathit{ic}}
\def\Picardfunctor{\mathrm{Pic}}
\def\Deformationcategory{\mathcal{D}\!\mathit{ef}}
\usepackage{cleveref}
\crefname{originallemma}{original lemma}{original lemmas}
\Crefname{originallemma}{Original lemma}{Original lemmas}
\makeatletter
\newlabel{lemma-describe-limits}{{lemma-describe-limits}{source}{}{sourcepointer-1}{}}
\newlabel{lemma-describe-limits@cref}{{[originallemma][0][]lemma-describe-limits}{[1][0][]source}{}{}{}}
\makeatother
\title{Constructible subsets of an inverse limit of spectral spaces are pullbacks of constructible subsets}
\author{The Stacks Project Authors}
\date{stacks-project:topology.tex:local-snapshot-20261012; GNU FDL 1.2 or later}
\begin{document}
\maketitle
\noindent\textit{Editorial note.} All delivered mathematical text is the author's own, copied from the pinned Stacks Project chapter source sources/stacks-topology/pinned-original-source.tex; no statement or proof step was rewritten, reordered or removed.\par\medskip
\noindent\textit{Editorial note.} Preamble provenance: the Stacks Project keeps one single shared preamble for every chapter; this card takes it directly from the staged sources/stacks-topology/preamble.tex as three verbatim spans (lines 8-12, 14-18 and 154-217), so no derived file was written and nothing was added or redefined. The wrapper class and the base packages are supplied by the assembler.\par\medskip
\noindent\textit{Editorial note.} Bibliography provenance: neither the counted statement nor its proof contains a citation (this lane rescans the whole span before writing the card), so the card carries no bibliography and no entry was reconstructed.\par\medskip
\noindent\textit{Editorial note.} Selection provenance: this card was opened under the file-based control contract (CONTROL.md), with the mechanical pre-card gate run on its statement and proof interval before the card was written.\par\medskip
\noindent\textit{Allowed standard prerequisites (editorial).}\begin{itemize}
\item The description of a basis of the limit topology by pullbacks of open subsets. Used inside the counted proof (source lines 5090-5100). Expresses the given constructible subset in terms of finitely many pullbacks of open sets.
\end{itemize}\medskip
\ifdefined\hypertarget\hypertarget{sourcepointer-1}{}\fi
\noindent\textit{Source reference.} Original lemma lemma-describe-limits: sources/stacks-topology/pinned-original-source.tex, source lines 2591--2591. Navigation pointer only: the counted text refers to the Stacks item lemma-describe-limits, which is an existing labeled result of this same chapter and is not part of this card. The pointer carries the target's own label as its number and its definition line as its span; its content is not reproduced here.\par\medskip
\begin{lemma}
\label{lemma-inverse-limit-spectral-spaces-constructible}
Let $\mathcal{I}$ be a cofiltered category. Let $i \mapsto X_i$ be a diagram
of spectral spaces such that for $a : j \to i$ in $\mathcal{I}$
the corresponding map $f_a : X_j \to X_i$ is spectral. Let $X = \lim X_i$
with projections $p_i : X \to X_i$. Let $E \subset X$ be a constructible
subset. Then there exists an $i \in \Ob(\mathcal{I})$ and a constructible
subset $E_i \subset X_i$ such that $p_i^{-1}(E_i) = E$. If $E$ is open,
resp.\ closed, we may choose $E_i$ open, resp.\ closed.
\end{lemma}

\begin{proof}
Assume $E$ is a quasi-compact open of $X$. By
Lemma \ref{lemma-describe-limits} we can write $E = p_i^{-1}(U_i)$
for some $i$ and some open $U_i \subset X_i$.
Write $U_i = \bigcup U_{i, \alpha}$ as a union of quasi-compact opens.
As $E$ is quasi-compact we can find $\alpha_1, \ldots, \alpha_n$
such that $E = p_i^{-1}(U_{i, \alpha_1} \cup \ldots \cup U_{i, \alpha_n})$.
Hence $E_i = U_{i, \alpha_1} \cup \ldots \cup U_{i, \alpha_n}$ works.

\medskip\noindent
Assume $E$ is a constructible closed subset. Then $E^c$ is quasi-compact
open. So $E^c = p_i^{-1}(F_i)$ for some $i$ and
quasi-compact open $F_i \subset X_i$ by the result of the previous paragraph.
Then $E = p_i^{-1}(F_i^c)$ as desired.

\medskip\noindent
If $E$ is general we can write $E = \bigcup_{l = 1, \ldots, n} U_l \cap Z_l$
with $U_l$ constructible open and $Z_l$ constructible closed. By the result
of the previous paragraphs we may write $U_l = p_{i_l}^{-1}(U_{l, i_l})$
and $Z_l = p_{j_l}^{-1}(Z_{l, j_l})$
with $U_{l, i_l} \subset X_{i_l}$ constructible open and
$Z_{l, j_l} \subset X_{j_l}$ constructible closed. As $\mathcal{I}$ is
cofiltered we may choose an object $k$ of $\mathcal{I}$ and morphism
$a_l : k \to i_l$ and $b_l : k \to j_l$. Then taking
$E_k = \bigcup_{l = 1, \ldots, n}
f_{a_l}^{-1}(U_{l, i_l}) \cap f_{b_l}^{-1}(Z_{l, j_l})$
we obtain a constructible subset of $X_k$ whose inverse image in $X$ is $E$.
\end{proof}

\end{document}
